\documentclass[10pt,aspectratio=169,serif]{beamer}

\usepackage[fontset=none, scheme=plain]{ctex}
\IfFontExistsTF{Songti SC}{%
	\setCJKmainfont{Songti SC}[AutoFakeSlant=0.2]%
	\setCJKsansfont{Songti SC}[AutoFakeSlant=0.2]%
}{%
	\setCJKmainfont{FandolSong-Regular}[Extension=.otf,BoldFont=FandolSong-Bold,AutoFakeSlant=0.2]%
	\setCJKsansfont{FandolSong-Regular}[Extension=.otf,BoldFont=FandolSong-Bold]%
}
\usepackage{amsmath,amssymb,graphicx,tikz,booktabs,multicol,calligra}
\usepackage{pgfplots}
\pgfplotsset{compat=1.18}
\usepackage{tcolorbox}
\tcbuselibrary{skins}
\usepackage{XyygU}
\usepackage{caption}
\captionsetup[figure]{name=图,labelsep=period}

\usetikzlibrary{arrows.meta,positioning,calc}

% 公式间距
\setlength{\abovedisplayskip}{0pt plus 2pt minus 2pt}
\setlength{\belowdisplayskip}{0pt plus 2pt minus 2pt}

\title[2.1 等式性质与不等式性质·练习]{2.1 等式性质与不等式性质（练习）}
\subtitle{不等关系与比较大小 \& 不等式性质}
\author{}
\institute{}
\date{2026-10-16}

\begin{document}

\begin{frame}
  \titlepage
\end{frame}

% ================= 第一组练习（教材P40） =================
\begin{frame}[t]{练习一：用不等式表示不等关系}
\textbf{1.} 用不等式或不等式组表示下面的不等关系：

\vspace{0.5em}
(1) 某高速公路规定通过车辆的车货总高度 $h$（单位：m）从地面算起不能超过 $4$ m；

(2) $a$ 与 $b$ 的和是非负实数；

(3) 如图，在一个面积小于 $350\ \mathrm{m^2}$ 的矩形地基的中心位置上建造一个仓库，仓库的四周建成绿地，仓库的长 $L$（单位：m）大于宽 $W$（单位：m）的 $4$ 倍。

\vspace{0.3em}
\begin{center}
\begin{tikzpicture}[scale=0.7]
  \draw[thick,fill=green!20] (0,0) rectangle (8,4);
  \draw[thick,fill=brown!40] (1.5,1) rectangle (6.5,3);
  \node at (4,2) {仓库};
  \node[left] at (0.75,2) {绿地};
  \draw[<->] (1.5,-0.3) -- (6.5,-0.3);
  \node[below] at (4,-0.3) {$L$};
  \draw[<->] (-0.3,1) -- (-0.3,3);
  \node[left] at (-0.3,2) {$W$};
  \node[above] at (4,4) {$5\,\mathrm{m}$};
  \draw[<->] (6.7,0) -- (6.7,1);
  \node[right] at (6.7,0.5) {$5\,\mathrm{m}$};
\end{tikzpicture}
\end{center}
\end{frame}

\begin{frame}[t]{练习一 第1题 解答}
\begin{tcolorbox}[colback=orange!10,colframe=orange,title=\textbf{解}]
(1) 车货总高度不能超过 $4$ m，即 $h$ 不大于 $4$：
\[
\alert{0 \le h \le 4}.
\]

\vspace{0.5em}
(2) $a$ 与 $b$ 的和是非负实数：
\[
\alert{a + b \ge 0}.
\]

\vspace{0.5em}
(3) 仓库长为 $L$、宽为 $W$，四周绿地宽均为 $5$ m，故矩形地基的长为 $L+10$，宽为 $W+10$。由题意：
\[
\begin{cases}
(L+10)(W+10) < 350,\\[2pt]
L > 4W,\\[2pt]
L > 0,\ W > 0.
\end{cases}
\]
\end{tcolorbox}
\end{frame}

% ========== 练习一 第2题 ==========
\begin{frame}[t]{练习一 第2题：作差比较大小}
\textbf{2.} 比较 $(x+3)(x+7)$ 和 $(x+4)(x+6)$ 的大小。
\end{frame}

\begin{frame}[t]{练习一 第2题 解答}
\begin{tcolorbox}[colback=orange!10,colframe=orange,title=\textbf{解：作差法}]
作差：
\begin{align*}
&(x+3)(x+7) - (x+4)(x+6)\\
&= (x^2 + 10x + 21) - (x^2 + 10x + 24)\\
&= -3 < 0.
\end{align*}

\vspace{0.5em}
因为差恒为 $-3 < 0$，所以
\[
\alert{(x+3)(x+7) < (x+4)(x+6)}.
\]
\small（与 $x$ 的取值无关）
\end{tcolorbox}
\end{frame}

% ========== 练习一 第3题 ==========
\begin{frame}[t]{练习一 第3题：证明平均值不等式}
\textbf{3.} 已知 $a>b$，证明：
\[
a > \frac{a+b}{2} > b.
\]
\end{frame}

\begin{frame}[t]{练习一 第3题 解答}
\begin{tcolorbox}[colback=orange!10,colframe=orange,title=\textbf{证明}]
\textbf{左边不等式：} 作差
\[
a - \frac{a+b}{2} = \frac{2a - a - b}{2} = \frac{a-b}{2}.
\]
因为 $a>b$，所以 $a-b>0$，从而 $\dfrac{a-b}{2}>0$，即 $a>\dfrac{a+b}{2}$。

\vspace{0.5em}
\textbf{右边不等式：} 作差
\[
\frac{a+b}{2} - b = \frac{a+b - 2b}{2} = \frac{a-b}{2} > 0.
\]
即 $\dfrac{a+b}{2} > b$。

\vspace{0.5em}
综上：$\alert{a > \dfrac{a+b}{2} > b}$。
\small（即算术平均数介于两数之间）
\end{tcolorbox}
\end{frame}

% ================= 第二组练习（教材P42） =================
\begin{frame}[t]{练习二：证明不等式的性质}
\textbf{1.} 证明不等式的性质 $1,\ 3,\ 4,\ 6$。

\vspace{0.8em}
\small
\textbf{性质1（对称性）：} $a>b \Rightarrow b<a$；

\textbf{性质3（可加性）：} $a>b \Rightarrow a+c>b+c$；

\textbf{性质4（可乘性）：} $a>b,\ c>0 \Rightarrow ac>bc$；\quad $a>b,\ c<0 \Rightarrow ac<bc$；

\textbf{性质6（同向同正可乘性）：} $a>b>0,\ c>d>0 \Rightarrow ac>bd$。
\end{frame}

\begin{frame}[t]{练习二 第1题 解答}
\begin{tcolorbox}[colback=orange!10,colframe=orange,title=\textbf{证明}]
\textbf{性质1：} 由 $a>b$，根据两实数大小关系的基本事实，$a-b>0$。而 $b-a=-(a-b)<0$，所以 $b<a$。

\vspace{0.3em}
\textbf{性质3：} $(a+c)-(b+c)=a-b>0$，所以 $a+c>b+c$。

\vspace{0.3em}
\textbf{性质4：} $ac-bc=(a-b)c$。当 $c>0$ 时，$(a-b)c>0$，故 $ac>bc$；当 $c<0$ 时，$(a-b)c<0$，故 $ac<bc$。

\vspace{0.3em}
\textbf{性质6：} 由 $a>b>0$ 及 $c>0$，据性质4得 $ac>bc$；由 $c>d>0$ 及 $b>0$，据性质4得 $bc>bd$。再由性质2（传递性）：$ac>bc>bd$，即 $ac>bd$。
\end{tcolorbox}
\end{frame}

% ========== 练习二 第2题 ==========
\begin{frame}[t]{练习二 第2题：填空}
\textbf{2.} 用不等号 $">"\text{ 或 }"<"$ 填空：

\vspace{0.8em}
(1) 如果 $a>b$，$c<d$，那么 $a-c$ \underline{\hspace{1.5cm}} $b-d$；

\vspace{0.5em}
(2) 如果 $a>b>0$，$c<d<0$，那么 $ac$ \underline{\hspace{1.5cm}} $bd$；

\vspace{0.5em}
(3) 如果 $a>b>0$，那么 $\dfrac{1}{a^2}$ \underline{\hspace{1.5cm}} $\dfrac{1}{b^2}$；

\vspace{0.5em}
(4) 如果 $a>b>c>0$，那么 $\dfrac{c}{a}$ \underline{\hspace{1.5cm}} $\dfrac{c}{b}$。
\end{frame}

\begin{frame}[t]{练习二 第2题 解答}
\begin{tcolorbox}[colback=orange!10,colframe=orange,title=\textbf{解}]
(1) 由 $c<d$ 得 $-c>-d$；又 $a>b$，据性质5（同向相加）：
$a+(-c)>b+(-d)$，即 $\alert{a-c>b-d}$。

\vspace{0.4em}
(2) 由 $c<d<0$ 得 $-c>-d>0$。又 $a>b>0$，据性质6：$a(-c)>b(-d)$，即 $-ac>-bd$，两边乘 $-1$ 反向：$\alert{ac<bd}$。

\vspace{0.4em}
(3) 由 $a>b>0$ 得 $a^2>b^2>0$（性质7），两边取倒数：$\dfrac{1}{a^2}<\dfrac{1}{b^2}$，即 $\alert{\dfrac{1}{a^2}<\dfrac{1}{b^2}}$。

\vspace{0.4em}
(4) 由 $a>b>0$ 得 $\dfrac{1}{a}<\dfrac{1}{b}$；又 $c>0$，两边乘 $c$：$\dfrac{c}{a}<\dfrac{c}{b}$，即 $\alert{\dfrac{c}{a}<\dfrac{c}{b}}$。
\end{tcolorbox}
\end{frame}

\begin{frame}[c]
  \begin{center}
    {\Huge\calligra Thanks!}
  \end{center}
\end{frame}

\end{document}
